\section{The full protocol, which text governs, and reporting} \label{sec:protocol}

\runin{Where the rules are.} Sections~\ref{sec:question} to \ref{sec:protocol} are the main part. Appendices~\ref{app:data}
to \ref{app:codemap} state every rule in full and are part of the pre-registration. The complete pre-registration, with
the appendices, is \code{PREREG\_S110\_FINAL.pdf}; a separate reader-oriented overview is issued as
\code{PREREG\_S110\_SHORT.pdf}. It introduces no additional rules. Appendices A--I give data, processing, thresholds, variants, localisation, procedure, public checks,
fixed wording and commands; J--N give dry runs, limits, history, files and the code's section map.

\runin{Who computes the outcome, and which text governs.} Frozen scripts (deposited with checksums and not changed
afterwards) compute every outcome from valid inputs: \code{s110\_final.py} runs the photometric test,
\code{s110\_localise.py} the pixel localisation, and \code{s110\_summary.py} writes the only sentences allowed to
describe the outcomes. Three rules of precedence apply.
\begin{enumerate}
\item If this text and the frozen code disagree on how an outcome is computed from valid inputs, the code defines the
  outcome. This text defines which runs and inputs are valid (Sections~\ref{sec:ephem} and \ref{sec:times} for the
  ephemeris and predicted times, Section~\ref{sec:stat} for the statistic, and Appendices~\ref{app:data}, \ref{app:processing}, \ref{app:variants},
  \ref{app:localisation}, \ref{app:procedure}, \ref{app:public} and \ref{app:commands}). The code checks only part of
  these rules; the rest are checked through the public record (Appendix~\ref{app:public}), and a few outcomes and
  statements are written by hand from posted files (Appendices~\ref{app:p1}, \ref{app:p5} and \ref{app:p7}). A result obtained in violation of any pre-registered rule
  is invalid and cannot be used to accept or reject either signal.
\item Where Sections~\ref{sec:question} to \ref{sec:protocol} and an appendix differ in any way, or a summary leaves
  something out, the formal rules of the appendix govern.
\item This version 1.1 is the authoritative text of the pre-registration. It replaces the earlier text of 2026 October
  5 (below).
\end{enumerate}

\runin{Order, failures and reporting.} No Sector 110 product is opened before the primary run. The first valid
primary run that writes an output defines the result; it and its checksum are posted before variants or localisation.
Applicable variants run in one continuous session in order (a), (b), (c), (e), (f), then (d) if public; otherwise (d) waits for the
QLP release public on 2027 February 28. Each attempt has an unedited record with command, input and script hashes,
UTC times, exit status and full output. A successful output cannot be replaced by a failure (Appendices F and G).

Before each run, the time-system headers are recorded and checked: TDB, SOLARSYSTEM, reference 2457000, fractional
reference 0 or absent, unit d, and cadence within 1\,\% of 2\,min (SPOC 2-min light curve or pixel file) or 200\,s (QLP/TESS-SPOC). A mismatch gives O4
for the primary, an unavailable variant, or localisation not possible, without conversion or product switching.
An environment or filename refusal is corrected and does not count as an attempt that failed while computing.
A computing failure or interruption before output is repeated once unchanged; a second failure gives O4 (primary), unavailable
(variant), or localisation not possible. A reading adaptation may change only how the file is read; its diff and
exact archival reproduction check must be posted before use. Helpers and individual variants are never adapted.
The binding reading limit is 30 days from release of the affected file, ending at 24:00 UTC; without a qualifying
adaptation it gives O4 (primary) or localisation not possible. A missing pixel file becomes final as not possible on or after 2027 February 28, with a saved search
showing none released by that deadline. The full cases and fixed failure sentences are in Appendix G.

The 30-day primary-run, 7-day posting and 14-day variant/localisation limits are administrative: report delays.
Order, the 2027 February 28 conditions and reading limits are binding. Invalid results are posted as invalid and
cannot accept or reject a signal; if a valid repeat is no longer allowed, the affected outcome is O4 (or localisation
not possible), reason ``pre-registered procedure not followed''. All outcomes, untested transits, variants and
localisations are reported with the fixed wording, including the variant-outcome line wherever ``robust'' is quoted.
The deposit is unchanged after publication; later ideas are separately labelled analyses (Appendices F--H).

\runin{Record status.} Version 1.1 replaces the earlier text deposited on 2026 October 5
(\href{https://doi.org/10.5281/zenodo.23171456}{doi:10.5281/zenodo.23171456}) and deleted the same day; that DOI retains a tombstone.
The public timestamp of this pre-registration is the date of the present deposit.
Unpublished internal drafts v1 and v2 are not part of this deposit and define no outcome.
The scientific choices made after examining archival data are documented in Appendix~\ref{app:history};
editorial changes and the detailed chronology are in \code{CHANGELOG.md}. Frozen code, predictions,
calibrations and archival outputs are unchanged in v1.1.
